\documentclass{standalone}

\usepackage{circuitikz}
\newcommand{\voltmeter}{node[circle,draw]{{\color{black!70}\textbf{V}}}}
\newcommand{\amperemeter}{node[circle,draw]{{\color{black!70}\textbf{A}}}}
\tikzset{
	pics/interruptor/.style={
		code={
			\draw (0,0) -- (1,0);
			\draw (2,0) -- (3,0);
			\draw (1,0) -- (1.8,0.5);
			\fill (1,0) circle (2pt);
			\fill (2,0) circle (2pt);}}}
\begin{document}
%% ---------------------------------------------------------------------------
		\begin{circuitikz}
		% Pile et interrupteur
		\draw [line width=1] (-3,0) to[battery, l=$Pile$, i=$I_g$, line width=0.5] (7,0);
		\pic[rotate=270, line width=1] at (7,0) {interruptor};
		\node at (7.7,0.5) {$S_1$};

		% Branche du bas avec Rh, R et ampèremètre
		\draw [line width=1] (7,0) --(7,-1);%-- (7,-4)
		\draw [line width=1] (7,-3)--(7,-7)--(5,-7)
		to[variable resistor ,l_=$Rh$, line width = 0.5, invert] (4,-7)
		to[ resistor, l_=$R$, line width = 0.5, invert] (-1.1,-7);
		\draw[line width=1] (-1.5,-7) \amperemeter (-1,-7);
		\draw[line width = 1] (-1.9,-7)--(-3,-7)--(-3,0);

		% Voltmètre en parallèle sur R
		\draw [line width=1] (5,0) -- (5,-3) -- (3.4,-3);
		\draw [line width=1] (3,-3) \voltmeter (2.63,-3) -- (-3,-3);

		% Flèches de tension
		\draw[>=stealth,->, thick] (-2,1.2) -- (6,1.2) node[midway,above] {$U_{g}$};
		\draw[>=stealth,->, thick] (-2.4,-6) -- (2.8,-6) node[midway,above] {$U_{R}$};
		\draw[>=stealth,->, thick] (3.3,-5.8) -- (6.3,-5.8) node[midway,above] {$U_{Rh}$};
	\end{circuitikz}

%% ---------------------------------------------------------------------------
\end{document}

